\section{Fundamentos da tomografia de tempo de trânsito}
\label{sec_fttt}
Como explicado anteriormente, \textit{Travel-Time Tomography} (TTT) é uma técnica que utiliza o atraso --- tempo de viagem --- entre a produção da onda na fonte e a sua chegada nos receptores, para estimar o caminho percorrido, e disso a velocidade da onda ao longo dessa rota.

\begin{figure}[H]
	\centering
	\begin{tikzpicture}[scale=1.5]
		\coordinate (tl) at (-3, 2);
		\coordinate (tr) at (3, 2);
		\coordinate (bl) at (-3, -2);
		\coordinate (br) at (3, -2);
		\coordinate (cc) at (0,0);
		\coordinate (tc) at (0,2);
		\coordinate (bc) at (0,-2);
		\coordinate (cl) at (-3, 0);
		\coordinate (cr) at (3,0);
		\coordinate (source) at ($(bl)!0.95!(br)$);
		\coordinate (receiver) at ($(bl)!0.9!(tl)$);
		\draw[gray,fill=none] (tl) rectangle (br);
		\draw[->-,thick] (source) to[out=120,in=-18] (receiver) node[pos=0.5,anchor=center,black] {$\bvr[v]$};
		%		\filldraw[blue!50!black] (receiver) rectangle (2,2);
		%		\filldraw[green!50!black] (source) circle (3pt);
		\node[square,fill=blue!70!black] () at (receiver) {};
		\node[triangle,fill=green!70!black] () at (source) {};
		
	\end{tikzpicture}
	\caption{Wave path on heterogeneous and continuously variable environment.}
	\label{fig:sec2:wavepath_heterogeneous-variable_environment}
\end{figure}

Consideremos uma onda que siga o caminho da \cref{fig:sec2:wavepath_heterogeneous-variable_environment}, representado pela curva $\bvr[v]$. Assumindo que a velocidade do som $c(\bvp)$ seja uma função do espaço ($\bvp = [x,y,z]$), então o tempo de viagem da onda que percorre o caminho $\bvr[v]$ é
\begin{equation}
	t_{\bvr[v]} = \int_{\bvr[v]} \frac{1}{c(\bvp)} \dd\bvp.
\end{equation}

O objetivo é, dadas $N$ medições de tempo de viagem $t_{\bvr[v]{i}}$ ($0 \leq 1 < N$ representando diferentes caminhos), estimar o campo de velocidade do som (SSF) $c(\bvp)$. Diferentes métodos podem ser utilizados para resolver esse problema, incluindo aproximação por funções base, por curvas Spline, discretização espacial, entre outras. Essas técnicas também podem ser aplicadas em situações 1D até 4D, e podem contar com diversidade espacial através de tanto múltiplas fontes quanto múltiplos receptores --- ou ambos.

Em específico, a técnica de TTT se baseia em minimizar o erro entre o tempo de viagem real medido, e a estimativa do tempo de viagem utilizando o SSF reconstruído --- esse conhecido como \textsc{problema inverso}. Para obter essa estimativa de tempo de viagem, é necessário o uso de um solver para encontrar os caminhos percorridos entre os dispositivos\footnote{Esse chamado de um solver da equação Eikonal, ou um solver Eikonal.}, processo chamado de raytracing --- esse conhecido como \textsc{problema direto}. Há diferentes abordagens que podem ser tomadas para construção do problema matemático basilar e que se baseiam exclusivamente (ou primariamente) informações de tempo de trânsito, e que busquem realizar a tomografia (reconstrução do SSF) do ambiente medido; contudo, o nome de \textit{Travel-Time Tomography} é reservado a técnicas como descritas anteriormente.

A maioria dos modelos (tanto TTT quanto técnicas adjacentes) focam em estimar o campo de vagarosidade (ou lentidão; o inverso da velocidade). Como a vagarosidade e a velocidade são quantidades recíprocas, a sigla SSF será utilizada indiscriminadamente para descrever ambos.

\subsection{Revisão de artigos sobre TTT}
\label{sec_ttt:subsec_revisao}
Como determinado anteriormente, a técnica de TTT é versátil devido à sua simplicidade, sendo aplicável em uma grande gama de aplicações, de sismologia acústica a tomografia computadorizada de tecidos, e avaliação eletromagnética de propriedades da água. 

\begin{itemize}
	\citem{aki_determination_1976,aki_determination_1977} - Esses dois trabalhos são considerados os artigos pioneiros da TTT, estudando classificação de camadas geológicas através da estimação do SSF, e estimação do epicentro de terremotos.
	\citem{ali_opensource_2019} - Procura estimar o SSF em tecidos utilizando tomografia computadorizada com ultrassom, utilizando \textit{Fast Marching Method} (FMM) para raytracing e como solver Eikonal.
	\begin{itemize}
		\citem{hormati_robust_2010} - Trata de um problema semelhante, utilizando diferentes abordagens para resolver o problema inverso.
		\citem{jeong_investigating_2024} - Compara TTT com tomografia por reflexão, usando ambas para treinamento de uma CNN.
		\citem{hellmann_borehole_2023} - Usa TTT para estimação do SSF em furos de sondagem em geleiras, com um passo de correção da posição dos dispositivos ao longo do furo.
	\end{itemize}
	\citem{dines_computerized_1979} - Trabalha com TTT, mas estima o campo de atenuação junto do SSF. Trabalha com sinais eletromagnéticos, mas a teoria e a matemática são as mesmas. Também propõe uma solução iterativa em pares de dispositivos.
	\citem{phillips_traveltime_1991} - Estuda diferentes métodos e modificações na TTT, comparando seus resultados em métricas como rugosidade, erro, espalhamento; com cada método trocando performance global por melhoria em alguma métrica.
	\citem{tang_travel_2024} - Propõe o uso da técnica \textit{Fast Sweeping Method} (FSM) ao invés de FMM para solução da equação Eikonal no problema direto, e foca primariamente na comparação dos dois solvers. Modelagem de reflexões com o FSM também é discutida no artigo.
	\citem{zhang_inversion_2016} - Propõe outra adaptação da TTT para inclusão de reflexões. O modelo minimiza uma combinação de vagarosidades média e aparente, além do tempo de regularização. Parte do princípio que, em certas aplicações, ondas refletidas podem chegar em momentos próximos que as diretas, permitindo seu uso.
\end{itemize}

\subsubsection{Temas em comum}
%
\def\X{\ensuremath{\mathsf{X}}}
%
A \cref{table:fttt_revisao_temas} compila uma lista dos temas em comum entre os artigos apresentados, e um pequeno comentário dos aspectos ou considerações que mais destacam cada.

\begin{table}[H]\small\centering
	\caption{Tabela comparativa dos artigos analisados. CG - Gradiente Conjugado (para minimização); $P_{\alpha}$ - Regularização; $\bvB{i}$ - Suavização do SSF.}
	\label{table:fttt_revisao_temas}
	\begin{tabular}{l c c c c c l}
		Paper 								& Discreto 	& Ray bending 	& CG 	& $P_{\alpha}$ 	& $\bvB{i}$ & Extra 	\\
		\cite{aki_determination_1976} 		& \X 		& -- 			& --	& \X			& --		& --		\\
		\cite{aki_determination_1977} 		& \X 		& -- 			& --	& \X			& --		& Análise de perturbação		\\
		\cite{ali_opensource_2019}			& \X		& \X			& \X	& \X			& \X		& Linearização estatística \\
		\cite{hormati_robust_2010}			& \X 		& \X			& \X	& \X			& --		& --		\\
		\cite{jeong_investigating_2024}	& \X		& \X			& --	& \X			& --		& --		\\
		\cite{hellmann_borehole_2023}		& \X		& \X			& --	& \X			& --		& Estimação de posição	\\
		\cite{dines_computerized_1979}		& \X		& --			& --	& --			& --		& Considera atenuação\\
		\cite{phillips_traveltime_1991}	& \X		& --			& --	& \X			& \X		& Compara várias técnicas	\\
		\cite{tang_travel_2024}			& \X		& \X			& \X	& --			& --		& Usa FSM \\
		\cite{zhang_inversion_2016}		& \X		& \X			& \X	& \X			& --		& Considera reflexões			
	\end{tabular}
\end{table}

\subsection{Método padrão para TTT discreto}
\label{sec_fttt:subsec_TTT-padrao}

Aqui será descrito um método padrão para a TTT, seguindo os artigos citados como guia. Diferenças em relação aos artigos podem surgir, como ignorar efeitos de \textit{ray bending}, tipo de regularização empregada, quantidade e posição dos dispositivos, propósito da técnica, ou outras particularidades.

Nesses quesitos, os artigos mais completos são \cite{ali_opensource_2019,phillips_traveltime_1991,hormati_robust_2010}, onde eles consideram regularização, \textit{ray bending}, e uma grande quantidade de dispositivos. Agora será estabelecido o terreno comum entre os artigos apresentados, as considerações medianas da modelagem, e algumas distinções importantes que são compartilhados entre alguns deles. A estrutura apresentada aqui não será idêntica a nenhum dos artigos citados, e busca destacar as similaridades e o ponto em comum deles.

\subsection{Estrutura básica}
Assumamos que o ambiente modelado é divido em uma grade de células $\sz{N_1}{N_2}$, com um total de $N = N_1 \cdot N_2$ células. Entre todas as fontes e receptores (ou entre algumas), um total de $M$ raios são produzidos e mensurados. Denota-se $\bvs \in \re^{\sz{N}{1}}$ como um vetor positivo representando a vagarosidade em cada célula, e $\bvr{m}(\bvs) \in \re^{\sz{N}{1}}$ um vetor positivo, o $i$-ésimo valor de $\bvr{m}(\bvs)$ representa a distância percorrida pelo $m$-ésimo raio na $i$-ésima célula. O vetor $\bvr{m}(\bvs)$ é geralmente obtido externamente através de algum solver Eikonal, considerando o princípio de Fermat de tempo mínimo, e a lei de Snell; ou obtido quase diretamente se for considerado um SSF com perturbações pequenas, em que efeitos de \textit{ray bending} podem ser ignorados, e assume-se um trajeto direto entre fonte e receptor.

O tempo de viagem para o $m$-ésimo $t_m(\bvs)$ é dado por
\begin{equation}
	t_m(\bvs) = \bvr[T]{m}(\bvs) \bvs
\end{equation}

Note que tanto o vetor de distâncias $\bvr{m}(\bvs)$ e o tempo de viagem $t_m(\bvs)$ dependem do SSF $\bvs$, já que um campo diferente levaria a caminhos diferentes, e também a um tempo de trânsito diferente.

Agora nós consideramos todos os $M$ raios disponíveis, do qual define-se $\bvt(\bvs) \in \re^{\sz{M}{1}}$ como um vetor dos tempos de viagem para o modelo $\bvs$ atual; e $\bvR{\bvs} \in \re^{\sz{M}{N}}$ uma matrix, em que a sua $m$-ésima linha é $\bvr[T]{m}(\bvs)$. Nós também denotamos $\bvt{\obs}$ como o vetor observado de tempo de trânsito, obtido a partir dos raios medidos. Com essas variáveis, o problema de minimização basal torna-se minimizar a função custo $E(\bvs)$, isso traduzindo-se em
\begin{equation}
	\opt{\bvs} = \arg\min_{\bvs} E(\bvs) = \norm{\bvR(\bvs) \bvs - \bvt{\obs}}^2
	\label{eq:sec2:def_minimization_bvz}
\end{equation}
onde $\opt{\bvs}$ é a solução ótima para a minimização estabelecida.

Note que esse problema é não-linear e recursivo: $\bvR(\bvs)$ depende do campo $\bvs$, e o modelo do campo $\bvs$ depende do resultado obtido pelo solver Eikonal $\bvR(\bvs)$. Posto isso, a solução é geralmente buscada iterativamente. O problema inverso é resolvido através da minimização
\begin{subgather}{eqs:sec2:iterative_minimization}
			\bvs{i+1} = \arg\min_{\bvs{i+1}} E(\bvs{i+1}) = \norm{\bvR(\bvs{i}) \bvs{i+1} - \bvt{\obs}}^2 
			\label{subeq:sec2:iterative_minimization:eq1}\\
			\nabla E(\bvs{i+1}) = 2\tr{\bvR}(\bvs{i}) \pts{\bvR(\bvs{i}) \bvs{i+1} - \bvt{\obs}}
			\label{subeq:sec2:iterative_minimization:eq2}
\end{subgather}
em que, na $(i+1)$-ésima iteração, $\bvs{i}$ (e disso $\bvR(\bvs{i})$) é tratado como constante, e o gradiente é tomado em relação a $\bvs{i+1}$. O problema direto é resolvido com algum solver Eikonal externamente. Assumindo invertibilidade de $\bvR[T](\bvs{i})\bvR(\bvs{i})$, então a solução da minimização --- o que se traduz em $\nabla E(\bvs{i+1}) = \bv0$ --- é
\begin{subgather}{eqs:sec2:inversion_zi+1}
		\bvs{i+1} = \bvG(\bvs{i}) \bvt{\obs} \\
		\bvG(\bvs{i}) = \inv*{\tr{\bvR}(\bvs{i}) \bvR(\bvs{i})} \tr{\bvR}(\bvs{i})
\end{subgather}

Como $\bvR(\bvs{i})$ é atualizado repetitivamente, $\bvG(\bvs{i})$ também será, com uma solução iterativa. Embora $\bvR(\bvs{i})$ dependa diretamente do SSF estimado, a dependência no iterador $i$ é muito mais prevalente no problema de minimização. Por isso, por questão de notação, denotaremos $\bvR{i} \equiv \bvR(\bvs{i})$, e $\bvG{i} = \bvG(\bvs{i})$.

\subsubsection{Minimização com posto deficiente}
Para a \cref{eqs:sec2:inversion_zi+1}, é crucial que $\bvR[T]{i}\bvR{i}$ seja uma matriz inversível. Dado que $\bvR{i} \in \re^{\sz{M}{N}}$, então $\bvR[T]{i}\bvR{i} \in \re^{\sz{N}{N}}$; é trivial que $\rank{\bvR{i}} = N$ e $N \leq M$ são condições necessárias e suficientes para a inversão ser possível. Intuitivamente, isso traduz-se em haver mais medições $M$ do que variáveis $N$, que ao menos $N$ medições sejam significativas, e que cada célula seja representada (visitada) por ao menos um raio. Essas condições, embora factíveis, não são práticas ou por vezes verificáveis, e portanto não é seguro assumir que a inversão seja viável.

Nessas situações, algumas alternativas já foram apresentadas na literatura. As mais populares são a solução indireta com gradiente descendente \cite{ali_opensource_2019,hormati_robust_2010,zhang_inversion_2016,tang_travel_2024}, ou uso de funções de regularização \cite{ali_opensource_2019,hormati_robust_2010,tang_travel_2024,zhang_inversion_2016}. Outra opção, que ainda não foi explorada na literatura no contexto de TTT, é utilizando a decomposição em valores singulares (SVD) \cite{barata_moore_2012,golub_calculating_1965} para a inversão \cite{brand_fast_2006,chen_treatment_2006}.

A partir daqui, será assumido que o kernel de inversão (modificado com regularização ou outros métodos) é invertível; situações em que isso não se aplica devem ser tratadas com inversões indiretas.

\subsubsection{Adição de regularização}

Como citado, uma alternativa para tratar o problema mal definido da inversão é a adição de regularização, simultaneamente aumentando as chances de uma solução única ser encontrada, e melhorando alguma condição secundária, como suavidade, erro em comparação a um modelo-base, e outras métricas. Escolher uma regularização adequada também pode ajudar a generalizar a solução encontrada para o caso contínuo, separando o problema de minimização da discretização espacial escolhida \cite{zhang_nonlinear_1998,delprat-jannaud_illposed_1993}.

A partir de \cref{eq:sec2:def_minimization_bvz}, podemos adicionar um termo de regularização na forma de $P_{\alpha}(\bvs{i+1})$, de modo que a forma iterativa de $\bvs{i+1}$ de \cref{eqs:sec2:iterative_minimization} se torne
\begin{subgather}		{eqs:sec2:iterative_minimization_Ezi}
		\bvs{i+1} = \arg\min_{\bvs{i+1}} E(\bvs{i+1}) = \norm{\bvR{i} \bvs{i+1} - \bvt{\obs}}^2 + P_\alpha(\bvs{i+1})\\		
		\nabla E(\bvs{i+1}) = 2\bvR[T]{i} \pts{\bvR{i} \bvs{i+1} - \bvt{\obs}} + \nabla P_{\alpha}(\bvs{i+1})
\end{subgather}
com $\alpha$ sendo um parâmetro de controle. O modelo mais comum para a função de regularização é da forma $P_{\alpha}(\bvs{i+1}) = \alpha^2\norm{\bvD\bvs{i+1}}^2$, onde $\bvD$ é uma matriz quadrada; e disso, a solução é
\begin{equation}
	\bvG{i} = \inv*{\tr{\bvR{i}} \bvR{i} + \alpha^2 \tr{\bvD}\bvD} \bvR[T]{i}
	\label{eq:sec2:solution_regularization-problem_simple}
\end{equation}

Para que esta etapa de inversão seja possível, basta que $\bvD$ seja uma matriz de posto completo\footnote{\label{footnote:positive-definite_inversion}Por definição, $\bvR[T]{i} \bvR{i}$ é uma matriz semidefinida positiva, $\bvD$ é definida positiva se $\bvD$ for uma matriz quadrada de posto completo, e a soma de uma matriz semidefinida positiva com uma matriz definida positiva é sempre definida positiva; esta sendo uma condição suficiente para a invertibilidade.}.
	
A matriz de regularização $\bvD$ pode ser escolhida para melhorar uma série de métricas. Escolher $\bvD = \Id$ (identidade) \cite{phillips_traveltime_1991} resulta na minimização considerando também a raiz quadrada média (RMS) do vetor de lentidão; isso é equivalente a considerar a presença de ruído branco nas medições e trocar a precisão do modelo pela rejeição de ruído. Escolher $\bvD z \approx \nabla^2 s(x,y)$, aproximando (discretamente) o Laplaciano do SSF \cite{zhang_nonlinear_1998,ali_opensource_2019}, resulta em uma minimização que penaliza o Laplaciano do campo de vagarosidade obtido; isto é, uma solução que também considera a curvatura/saltos no SSF, buscando um campo mais suave.

\cite{zhang_nonlinear_1998} também comenta brevemente sobre o uso de derivadas de terceiro grau em vez de derivadas de segundo grau (sendo estas últimas a regularização Laplaciana), o que também poderia resultar em uma boa suavização do campo de velocidade. \cite{hormati_robust_2010} propõe uma função de regularização baseada na representação ortonormal da distribuição de lentidão e nas normas $L_0$ ou $L_1$.

Outra possibilidade que ainda está em aberto para exploração é o uso de múltiplas funções de regularização. Por exemplo, o uso combinado das duas funções apresentadas anteriormente (regularização de ruído branco e penalização Laplaciana) pode levar a uma boa suavidade na função de distribuição de lentidão, mantendo um kernel invertível garantido. Ou seja, dizemos que
\begin{equation}	
	E(\bvs{i+1}) = \norm{\bvR{i} \bvs{i+1} - \bvt{\obs}}^2 + \alpha^2\norm{\bvD\bvs{i+1}}^2 + \beta^2 \norm{\bvs{i+1}}^2
\end{equation}
para o qual
\begin{equation}	
	\bvG{i} = \inv*{\tr{\bvR{i}} \bvR{i} + \alpha^2 \tr{\bvD}\bvD + \beta^2 \Id} \bvR[T]{i}
\end{equation}
O termo entre colchetes é garantidamente definido positivo (\cref{footnote:positive-definite_inversion}), e, portanto, invertível com segurança.

\subsubsection{Média convolucional/suavização por desfoque}

Outra ferramenta interessante encontrada na literatura \cite{phillips_traveltime_1991,meyerholtz_convolutional_1989} é o uso de um desfoque no gradiente. Ou seja, definimos que
\begin{equation}	
	E(\bvs{i+1}) = \norm{\bvR[d]{i} \bvs[d]{i+1} - \bvt{\obs}}^2 + P_\alpha(\bvs[d]{i+1})
\end{equation}
onde $\bvR[d]{i} = \bvR{i}\bvB{i}$ e $\bvs[d]{i+1} = \inv{\bvB{i}} \bvs{i+1}$, e $\bvB{i}$ é a matriz de desfoque da $i$-ésima iteração (baseada em um desfoque Gaussiano, ou qualquer outro kernel de média, que diminui de raio a cada iteração). Neste cenário, a solução $\bvs{i+1}$ é
\begin{equation}	
	\bvG{i} = \bvB{i} \inv*{\tr{\bvB{i}} \tr{\bvR{i}} \bvR{i} \bvB{i} + \alpha^2 \tr{\bvD}\bvD} \tr{\bvB{i}} \bvR[T]{i}
\end{equation}

Para alcançar isso, a função de regularização precisa ser adaptada, de modo que $P_{\alpha}(\bvs{i+1}) = \alpha^2 \norm{\bvD \inv{\bvB{i}} \bvs{i+1}}^2$; isso significa que a regularização é aplicada ao sinal inversamente desfocado. Conforme mencionado, a matriz de desfoque $\bvB{i}$ depende do iterador $i$. Normalmente, o raio de desfoque é reduzido com o aumento de $i$, visando encontrar gradualmente uma solução mais refinada e precisa (menos desfocada).

Essa média convolucional (também chamada de alargamento de raios \cite{meyerholtz_convolutional_1989,phillips_traveltime_1991}) pode auxiliar na busca de soluções ao dispersar os raios e espalhar os caminhos percorridos para células vizinhas, lidando assim com o problema de algumas células não serem atravessadas por nenhum raio.

Em \cite{ali_opensource_2019}, uma estrutura de desfoque semelhante é proposta, mas esse desfoque é aplicado diretamente ao gradiente, de modo que
\begin{equation}	
	\nabla E(\bvs{i+1}) = 2 \bvB{k} \bvR[T]{i}\pts{\bvR{i} \bvs{i+1} - \bvt{\obs}} + 2 \alpha^2 \tr{\bvD}\bvD \bvs{i+1}
\end{equation}

Como no artigo citado a solução $\bvs{i+1}$ é encontrada pelo método do gradiente conjugado (iterativo), e não por inversão, a matriz de desfoque é variável dentro da busca iterativa do gradiente conjugado pela inversa (com iterador $k$), em vez do processo iterativo de resolução de Eikonal (iterador $i$). Assumindo invertibilidade, a solução pode ser obtida diretamente através de
\begin{equation}		
	\bvG{i} = \inv*{\bvB{i} \bvR[T]{i} \bvR{i} + \tr{\bvD}\bvD}\bvB{i} \bvR[T]{i}
\end{equation}
situação em que o desfoque pode ser atenuado em cada iteração de $i$, ou mesmo mantido constante.

Note que ambas as estruturas de suavização apresentadas requerem a presença de um termo de regularização (caso contrário, o termo inverso pode ser expandido e os termos $\bvB{i}$ serão cancelados); ou que o sistema seja suficientemente mal condicionado de forma que uma minimização direta não seja possível, caso em que o efeito de suavização não se cancela necessariamente com seu inverso.